
\section{Lutherdom and Church Polity}

Editorial article in “Folkebladet” for September 30, 1896.

It is a matter which ought to be known by all Lutherans, that the Lutheran Church, on account of its unhappy connection with the State, did not succeed in carrying through its principles with respect to church polity.

The principles may nevertheless be just as sound and good, even if one did not have the courage and strength to carry them through. It is therefore not out of place now, when The United Church has taken so decided a stand for corporate power and against the freedom of the congregation, to look a little at the Lutheran principles in this matter.

For it is by no means altogether indifferent whether, in constitutional questions, we pass over onto Roman ground, or remain standing on Lutheran. For it is certain enough that, if once one slides over in the question of the Church and the ecclesial congregation into Roman modes of thought, one will soon be compelled to follow after in all other things.

The principal question for the Lutheran Church in the sphere of polity is twofold: on the one side it concerns preserving the freedom of ecclesial conviction, and on the other side not losing ecclesial unity.

For the Lutheran Church is from the beginning a party of opposition, or a minority. So exceedingly small was this minority, that it was one man against the whole Roman Church. And this one man therefore maintained that he alone had the same right to hold his conviction as the whole outward Church together with all its bishops, popes, priests, and councils. And that, when the question was to be decided whether he had the truth on his side or the Church had the truth, then the Church was not judge in its own cause, but only the Word of God was the highest authority, the only authority to which both sides were unconditionally bound to submit.

It is this mighty word of freedom which is, of course, just as intolerable to the fleshly Church now as then. This is free conviction’s great protest against the tyranny of society or Church. For the sake of this cause martyrs suffered death by the thousands at the stake and on the scaffold.

This, then, is a first principle in the Lutheran Church: the right of the individual believer and of the individual congregation to hold its Christian conviction of conscience unimpaired by the power of the Church and its authorities. If this is a sinful principle of rebellion, then of course all Lutheranism is condemned; for in the beginning it had the whole so-called Christian Church against it. It had only the Word of God and the Spirit of God on its side.

But the Lutheran Church does not regard it as indifferent whether Christian people and Christian congregations each go their own way, in disagreement with one another or not. It does not abolish unity for freedom’s sake, just as little as it abolishes freedom for unity’s sake. At the same time, therefore, as Lutherans maintain that one man with the Gospel on his side has an unconditional right to stand against the whole Church with all its authorities, when they depart from the Gospel, at the same time they also maintain that the proclamation of the Gospel, one, whole, and uncorrupted, must be found in all congregations and be the common bond which binds them together.

Here we have the two principles which in all constitutional work within the Lutheran Church must be inviolably maintained: freedom of personal conviction on the one side, unity in the confession of the Gospel on the other side.

A Lutheran Free Church must therefore disregard all state-church distortions of the congregation’s nature and outward form; and, in order to be faithful and have firm ground beneath its feet, it is compelled to go all the way back to the New Testament and follow the congregation’s footsteps there.

And in the New Testament we find a congregational order which precisely joins together the two things that for the Lutheran Church are the overwhelmingly chief matter of all: freedom of personal conviction and unity in the knowledge of the Gospel. The congregations have no outward or human authority over them; but inwardly within the congregation, and outwardly among the congregations, they preserve the unity of the Spirit without any superintending authority on earth.

Must not we begin on the same way, when we see how altogether impossible it is for our associations, with their annual meetings and officeholders, to preserve the freedom which nevertheless is the very life-air of Christianity and of the Christian Church?

Georg Sverdrup
